\documentclass{MFM-hw} % You must have the MFM-hw.cls file in your project or working directory (i.e. the same directory as this document) 

\HWauthor{Wei Qi}{226010267@link.cuhk.edu.cn} % Put your name and CUHK-SZ e-mail on this line
\HWno{1} % Enter the number of the homework you are working on
\HWcourse{MFM 5130} % Enter the course department and number here
\date{Sep 18, 2026}
\usepackage{amsmath}
\usepackage{amsfonts}
\usepackage{amssymb}
\usepackage{amsthm}
\usepackage{bbm} % For \mathbbm{1}
% \usepackage{hyperref}

\newcommand{\VaR}{\operatorname{VaR}}
\newcommand{\ES}{\operatorname{ES}}
\newcommand{\real}{\mathbb{R}}
\newcommand{\normal}{\operatorname{N}}
\newcommand{\N}{\mathcal{N}}
\newcommand{\dd}{\,\mathrm{d}}
\newcommand{\E}{\mathbb{E}}
\newcommand{\Par}{\operatorname{Par}}
\renewcommand{\P}{\mathbb{P}}


\begin{document}
\maketitle
\HWproblem
$\VaR_t : L^0 \rightarrow \real $ and $\ES_t : L^1 \rightarrow \real,\; t \in (0,1)$ are defined as follows:
\begin{align*}
    & \VaR_t(X) \triangleq F_X^{\leftarrow}(t) = \inf \{x \in \real : \P(X \leq x) \geq t \},\\
    & \ES_t(X) \triangleq \frac{1}{1-t} \int_t^1 \VaR_p(X) \dd p.
\end{align*}
Show that for fixed $t \in (0,1)$, both $\VaR_t$ and $\ES_t$ satisfy monotonicity, cash-additivity and positive homogeneity.\\

\textbf{Proof of monotonicity:} For any $X_1, X_2$ that satisfy $X_1 \leq X_2$ almost surely, there is a null set $N$ with
$$
X_1(\omega) \le X_2(\omega) \quad \text{for all } \omega \in N^c .
$$

Fix any $x \in \mathbb{R}$. On $N^c$,
$$
X_2(\omega) \le x \;\Longrightarrow\; X_1(\omega) \le X_2(\omega) \le x ,
$$
so $\{X_2 \le x\} \cap N^c \subseteq \{X_1 \le x\}$. Since $N$  is a null set,
$$
 \mathbb{P}(X_1 \le x) \;\ge\; \mathbb{P}(X_2 \le x,\; \omega \in N^c) = \mathbb{P}(X_2 \le x).
$$

Because for every $x$:
$$
\mathbb{P}(X_2 \le x) \ge t \;\Longrightarrow\;\mathbb{P}(X_1 \le x) \ge \mathbb{P}(X_2 \le x) \ge t,
$$
hence
\[
    \{x \in \real : \P(X_2 \leq x) \geq t \} \subseteq \{x \in \real : \P(X_1 \leq x) \geq t \},
\]
thus by definition of infimum,
\[
    \inf \{x \in \real : \P(X_1 \leq x) \geq t \} \leq \inf \{x \in \real : \P(X_2 \leq x) \geq t \},
\]
which is
\[
    \VaR_t(X_1) \leq \VaR_t(X_2),
\]
showing the monotonicity of $\VaR_t$. 

So for any $p \in (t, 1)$,
\[
    \VaR_p(X_1) \leq \VaR_p(X_2),
\]
for any $t \in (0,1)$, hence
\[
    \frac{1}{1-t} \int_t^1 \VaR_p(X_1) \dd p \leq \frac{1}{1-t} \int_t^1 \VaR_p(X_2) \dd p,
\]
which is
\[
    \ES_t(X_1) \leq \ES_t(X_2),
\]
showing the monotonicity of $\ES_t$.\qed \\

\textbf{Proof of cash-additivity: }For any $l \in \real$, let $u = x - l$. By definition of infimum,
\begin{align*}
    \VaR_t(X+l) &= \inf \{ x \in \real : \P (X+l \leq x) \geq t\} = \inf \{ x \in \real : \P (X \leq x-l) \geq t\}\\
    & = \inf \{ u + l : \P (X \leq u) \geq t\} = \inf \{ u \in \real : \P (X \leq u) \geq t\} + l \\
    & = \VaR_t(X) + l,
\end{align*}
which shows the cash-additivity of $\VaR_t$. Hence
\begin{align*}
    \ES_t(X+l) &= \frac{1}{1-t} \int_t^1 \VaR_p(X+l) \dd p = \frac{1}{1-t} \int_t^1 (\VaR_p(X) + l) \dd p \\
    & = \frac{1}{1-t}\left[\int_t^1 \VaR_p(X)\dd p + (1-t)l\right] = \frac{1}{1-t}\int_t^1 \VaR_p(X) \dd p + l \\
    & = \ES_t(X)+l,
\end{align*}
which shows the cash-additivity of $\ES_t$.\qed\\

\textbf{Proof of positive homogeneity: }For any $\lambda > 0$, let $u = x/\lambda$. By definition of infimum,
\begin{align*}
    \VaR_t(\lambda X) &= \inf \{ x \in \real : \P (\lambda X \leq x) \geq t\} = \inf \{ x \in \real : \P (X \leq \frac{x}{\lambda}) \geq t\}\\
    & = \inf \{ \lambda u : \P (X \leq u) \geq t\} = \lambda \inf \{ u \in \real : \P (X \leq u) \geq t\} \\
    & = \lambda \VaR_t(X),
\end{align*}
which shows the positive homogeneity of $\VaR_t$. Hence
\begin{align*}
    \ES_t(\lambda X) &= \frac{1}{1-t} \int_t^1 \VaR_p(\lambda X) \dd p = \frac{1}{1-t} \int_t^1 \lambda \VaR_p(X)  \dd p \\ 
    & = \frac{\lambda}{1-t}\int_t^1 \VaR_p(X) \dd p  \\
    & = \lambda \ES_t(X),
\end{align*}
which shows the positive homogeneity of $\ES_t$. \qed

\HWproblem
Show for $\alpha \in (0,1)$,
\begin{align*}
    &\ES_{\alpha}(L) = \frac{\E((L - F_L^{\leftarrow}(\alpha))_+)}{1-\alpha} +F_L^{\leftarrow}(\alpha),\\
    &\ES_{\alpha}(L) = \frac{\E(L\mathbbm{1}_{\{L>F_L^{\leftarrow}(\alpha)\}})+F_L^{\leftarrow}(\alpha)\bigl(1-\alpha-\bar{F}_L(F_L^{\leftarrow}(\alpha))\bigr)}{1-\alpha}.
\end{align*}
Hence show that if $F_L$ is continuous at $F_L^{\leftarrow}(\alpha)$, then
\[
    \ES_{\alpha}(L) = \frac{\E(L\mathbbm{1}_{\{L>F_L^{\leftarrow}(\alpha)\}})}{1-\alpha}.
\] \\

\textbf{Lemma: }Let $Y \in \mathcal{L}_1$ be a random variable on $(\Omega, \mathcal{F},\P)$, with cumulative distribution function $F_Y(y)$. Let 
\[
    F_Y^{\leftarrow}(u) = \inf \{y \in \real : F_Y(y) \ge u\},
\]
then we have
\[
    F_Y(F_Y^{\leftarrow}(u)) \ge u,\quad F_Y(y) \ge u \Leftrightarrow y \ge F_Y^{\leftarrow}(u), \quad \E(Y) = \int_0^1 F_Y^{\leftarrow}(u) du.
\]

\textbf{Proof of Lemma: }By definition of infimum,
\[
    F_Y(y) \ge u \Rightarrow y \ge F_Y^{\leftarrow}(u).
\]
Since CDF $F_Y(y)$ is non-decreasing and right-continuous, $\{y \in \real :F_Y(y) \ge u\} = [F_Y^{\leftarrow}(u), +\infty)$, which derives that
\[
    F_Y(F_Y^{\leftarrow}(u)) \ge u,
\]
hence
\[
    y \ge F_Y^{\leftarrow}(u) \Rightarrow F_Y(y) \ge F_Y(F_Y^{\leftarrow}(u)) \ge u. 
\]

Now we get
\[
    F_Y(y) \ge u \Leftrightarrow y \ge F_Y^{\leftarrow}(u).
\]
Let $U$ be a random variable independent from $Y$, satisfying uniform distribution on interval $[0,1]$. Let $Y^{*}(U) = F_Y^{\leftarrow}(U)$. Because
\[
    \P(Y^* \le y) = \P(F_Y^{\leftarrow}(U) \le y) = \P(U \le F_Y(y)) = F_Y(y),  
\]
which means $Y^*$ and $Y$ have the same CDF, So
\[
    \E(Y) = \E(Y^*) = \int_0^1 Y^*(u)\dd u = \int_0^1 F_Y^{\leftarrow}(u)\dd u.
\]\qed

\textbf{Proof of original proposition: }Define $Y \in \mathcal{L}_1$ (the problem uses $\E(Y)$, which means $Y \in \mathcal{L}_1$ in default) as
$$
    Y = (L- F_L^{\leftarrow}(\alpha))_{+} = 
    \begin{cases}
        0, & L \le F_L^{\leftarrow}(\alpha) \\
        L - F_L^{\leftarrow}(\alpha), &L >F_L^{\leftarrow}(\alpha)
    \end{cases}
$$
Because $L$ has a CDF $F_L(x)$. Fix $\alpha$, the CDF of $L - F_L^{\leftarrow}(\alpha)$ should be $F_L(x + F_L^{\leftarrow}(\alpha))$. Combined with the expression of $Y$, it derives that the CDF of $Y$
\[
    F_Y(y) = 
    \begin{cases}
        0, & y < 0,\\
        F_L(y + F_L^{\leftarrow}(\alpha)), &y \ge 0.
    \end{cases}
\]

$F_Y(y)$ is non-decreasing and right-continuous, and $F_Y(0) = F_Y(0+) = F_L(F_L^{\leftarrow}(\alpha)), F_Y(0-) = 0$. Hence
\begin{align*}
    F_Y^{\leftarrow}(u) &= 
    \begin{cases}
        0, &0 \le u <  F_L(F_L^{\leftarrow}(\alpha)),\\
        \inf\{y \in \real:F_L(y + F_L^{\leftarrow}(\alpha)) \ge u\}, & F_L(F_L^{\leftarrow}(\alpha)) \le u \le 1,
    \end{cases}\\
    &=
    \begin{cases}
        0, \qquad\qquad\qquad\quad &0 \le u <  F_L(F_L^{\leftarrow}(\alpha)),\\
        F_L^{\leftarrow}(u) - F_L^{\leftarrow}(\alpha), \qquad\qquad\qquad\quad & F_L(F_L^{\leftarrow}(\alpha)) \le u \le 1.
    \end{cases}
\end{align*}
It can be noticed that when $\alpha \le u \le F_L(F_L^{\leftarrow}(\alpha))$, $F_L^{\leftarrow}(u) = F_L^{\leftarrow}(\alpha)$. The above formula can also be written as
\begin{align*}
    F_Y^{\leftarrow}(u) = 
    \begin{cases}
        0, &0 \le u <  \alpha,\\
        F_L^{\leftarrow}(u) - F_L^{\leftarrow}(\alpha), & \alpha \le u \le 1,
    \end{cases}\\
\end{align*}

According to the lemma, 
\[
    \E(Y) = \int_0^1 F_Y^{\leftarrow}(u) du,
\]
where
\[
    \text{LHS} = \E((L - F_L^{\leftarrow}(\alpha))_{+})
\]
\[
    \text{RHS} = \int_{\alpha}^{1} (F_L^{\leftarrow}(u) - F_L^{\leftarrow}(\alpha)) \dd u  = [\ES_{\alpha}(L) - F_L^{\leftarrow}(\alpha)](1-\alpha),
\]
hence
\[
    \E((L - F_L^{\leftarrow}(\alpha))_{+}) = [\ES_{\alpha}(L) - F_L^{\leftarrow}(\alpha)](1-\alpha),
\]
which is
\begin{align*}
    \ES_{\alpha}(L) &= \frac{\E((L - F_L^{\leftarrow}(\alpha))_+)}{1-\alpha} +F_L^{\leftarrow}(\alpha) \\
    &=\frac{\E((L - F_L^{\leftarrow}(\alpha)){\mathbbm{1}_{\{L>F_L^{\leftarrow}(\alpha)\}}})}{1-\alpha} +F_L^{\leftarrow}(\alpha) \\
    &=\frac{\E(L\mathbbm{1}_{L>F_L^{\leftarrow}(\alpha)})-F_L^{\leftarrow}(\alpha)\E(\mathbbm{1}_{L>F_L^{\leftarrow}(\alpha)})}{1-\alpha} +F_L^{\leftarrow}(\alpha) \\
    &=\frac{\E(L\mathbbm{1}_{L>F_L^{\leftarrow}(\alpha)})-F_L^{\leftarrow}(\alpha)\P(L>F_L^{\leftarrow}(\alpha))}{1-\alpha} +F_L^{\leftarrow}(\alpha) \\
    &=\frac{\E(L\mathbbm{1}_{L>F_L^{\leftarrow}(\alpha)})-F_L^{\leftarrow}(\alpha)\bar{F}_L(F_L^{\leftarrow}(\alpha))}{1-\alpha} +F_L^{\leftarrow}(\alpha)\\
    &=\frac{\E(L\mathbbm{1}_{\{L>F_L^{\leftarrow}(\alpha)\}})+F_L^{\leftarrow}(\alpha)\bigl(1-\alpha-\bar{F}_L(F_L^{\leftarrow}(\alpha))\bigr)}{1-\alpha}, 
\end{align*}
where
\[
    \bar{F}_L(u) = \P(L>u) = 1 - F_L(u).
\]

If $F_L$ is continuous at $F_L^{\leftarrow}(\alpha)$, then $F_L(F_L^{\leftarrow}(\alpha)) = \alpha$, and $\bar{F_L}(F_L^{\leftarrow}(\alpha)) = 1-\alpha$. Substitute it into the expression of $\ES_\alpha(L)$ and we get
\[
    \ES_{\alpha}(L) = \frac{\E(L\mathbbm{1}_{\{L>F_L^{\leftarrow}(\alpha)\}})}{1-\alpha}.
\]\qed

\HWproblem
Suppose the random variable $L \sim \normal(\mu,\sigma^2)$, where $\sigma > 0$. Show that
\[
    \VaR_{\alpha}(L) = \mu + \sigma \Phi^{-1}(\alpha),
\]
\[
    \ES_{\alpha}(L) = \mu + \sigma\frac{\varphi(\Phi^{-1}(\alpha))}{1-\alpha},
\]
and
\[
    \lim\limits_{\alpha \rightarrow 1^{-}} \frac{\ES_{\alpha}(L)}{\VaR_{\alpha}(L)} = 1,
\]
where $\Phi$ and $\varphi$ are the CDF and PDF of a standard normal distribution.\\

\textbf{Proof: }Let random variable $X = (L - \mu)/\sigma$, then $X \sim \normal(0,1)$ and $L = \mu + \sigma X$. By translation invariance and positive homogeneity of $\VaR$ and $\ES$, we have
\[
    \VaR_{\alpha}(L) = \VaR_{\alpha}(\mu + \sigma X) = \mu + \sigma \VaR_{\alpha}(X) = \mu + \sigma \Phi^{-1}(\alpha).  
\]
\[
    \ES_{\alpha}(L) = \ES_{\alpha}(\mu + \sigma X) = \mu + \sigma \ES_{\alpha}(X),
\]
where since $\Phi$ is continuous,
\begin{align*}
    \ES_{\alpha}(X) &= \frac{\E(X\mathbbm{1}_{\{X > \Phi^{-1}(\alpha)\}})}{1-\alpha} = \frac{\int_{\Phi^{-1}(\alpha)}^{+\infty} x \varphi(x) \dd x}{1-\alpha} \\
    &= \frac{1}{1-\alpha} \int_{\Phi^{-1}(\alpha)}^{+\infty} \frac{x}{\sqrt{2\pi}}e^{\frac{-x^2}{2}} \dd x = \frac{1}{1-\alpha}\left(-\frac{1}{\sqrt{2\pi}}e^{-\frac{x^2}{2}}\right) \Bigg|_{\Phi^{-1}(\alpha)}^{+\infty}\\
    &= \frac{\varphi(\Phi^{-1}(\alpha))}{1-\alpha},
\end{align*}
hence
\[
    \ES_{\alpha}(L) = \mu + \sigma\frac{\varphi(\Phi^{-1}(\alpha))}{1-\alpha}.
\]

Let $u = \Phi^{-1}(\alpha)$, thus $\alpha = \Phi(u)$. when $\alpha \rightarrow 1^{-}$, $u \rightarrow +\infty$. Then
\[
    \lim\limits_{\alpha \rightarrow 1^{-}} \VaR_{\alpha}(L) = \lim\limits_{u \rightarrow +\infty} (\mu + \sigma u) = +\infty,
\]
and
\[
    \lim\limits_{\alpha \rightarrow 1^{-}} \ES_{\alpha}(L) = \lim\limits_{u \rightarrow +\infty} (\mu + \sigma \frac{\varphi(u)}{1-\Phi(u)}) ,
\]
where by L'Hôpital's rule
\[
     \lim\limits_{u \rightarrow +\infty}\frac{\varphi(u)}{1-\Phi(u)} = \lim\limits_{u \rightarrow +\infty}\frac{-u\varphi(u)}{-\varphi(u)} = \lim\limits_{u \rightarrow +\infty} u = +\infty.
\]
So
\[
    \lim\limits_{\alpha \rightarrow 1^{-}} \VaR_{\alpha}(L) = +\infty, \quad\lim\limits_{\alpha \rightarrow 1^{-}} \ES_{\alpha}(L) = +\infty.
\]

By L'Hôpital's rule,
\[
    \lim\limits_{\alpha \rightarrow 1^{-}} \frac{\ES_{\alpha}(X)}{\VaR_{\alpha}(X)} = \lim\limits_{u \rightarrow + \infty} \frac{\varphi(u)}{u[1-\Phi(u)]} = \lim\limits_{u \rightarrow + \infty} \frac{-u\varphi(u)}{1-\Phi(u)-u\varphi(u)} = \lim\limits_{u \rightarrow + \infty} \frac{1}{1-\frac{1-\Phi(u)}{u\varphi(u)}} = 1,
\]
where the last step uses $\varphi(u)/(1-\Phi(u)) \sim u$, which gives $(1-\Phi(u))/(u\varphi(u)) \sim u^{-2} \rightarrow 0$. Hence
\[
    \lim\limits_{\alpha \rightarrow 1^{-}} \frac{\ES_{\alpha}(L)}{\VaR_{\alpha}(L)} = \lim\limits_{\alpha \rightarrow 1^{-}}\frac{\mu + \sigma\ES_{\alpha}(X)}{\mu + \sigma \VaR_{\alpha}(X)} = \lim\limits_{\alpha \rightarrow 1^{-}}\frac{\mu/\ES_{\alpha}(X) + \sigma}{\mu/\VaR_{\alpha}(X) + \sigma} = 1.
\]\qed

\HWproblem
If $L \sim t_{\nu}$ (Student-t on $\nu$ degrees of freedom), that is
\[
    f_L(x) = \frac{\Gamma(\frac{\nu+1}{2})}{\sqrt{\nu\pi}\Gamma(\frac{\nu}{2})}\left(1+\frac{x^2}{\nu}\right)^{-\frac{\nu+1}{2}},\quad x \in \real, \nu > 0,
\]
show
\[
    1-F_L(x) \sim cx^{-\nu}, \quad x \rightarrow +\infty.
\]
Then show that for $\nu > 1$,
\[
    \lim\limits_{\alpha \rightarrow 1^{-}} \frac{\ES_{\alpha}(L)}{\VaR_{\alpha}(L)} = \frac{\nu}{\nu-1} > 1.
\]\\

\textbf{Proof: }When $x \rightarrow + \infty$, $1-F_L(x) \rightarrow 0$ and $x^{-v} \rightarrow 0$. By L'Hôpital's rule,
\[
\lim\limits_{x \rightarrow +\infty} \frac{1-F_L(x)}{x^{-\nu}} = \lim\limits_{x \rightarrow +\infty} \frac{f_L(x)}{\nu x^{-\nu-1}} = \lim\limits_{x \rightarrow +\infty}\frac{\Gamma(\frac{\nu+1}{2})}{\nu\sqrt{\nu\pi}\Gamma(\frac{\nu}{2})} \left(\frac{1}{x^2} + \frac{1}{\nu} \right)^{-\frac{\nu+1}{2}} = c,
\]
where c is a constant as
\[
c = \frac{\Gamma(\frac{\nu+1}{2})}{\sqrt{\nu\pi}\Gamma(\frac{\nu}{2})} \nu^{\frac{\nu-1}{2}},
\]
hence
\[
    1-F_L(x) \sim cx^{-\nu}, \quad x \rightarrow +\infty.
\]

Because $F_L(x)$ is continuous and strictly increasing, $F_L^{\leftarrow}(t) = F_L^{-1}(t)$ for any $t \in (0,1)$. Let $u = \VaR_{\alpha}(L) = F_L^{-1}(\alpha)$. Then when $\alpha \rightarrow 1^{-}$, $\VaR_{\alpha}(L) \rightarrow +\infty$ and $\ES_{\alpha}(L) \rightarrow + \infty$. We have
\[
    \ES_{\alpha}(L) = \frac{1}{1-\alpha}\int_{\alpha}^1F_L^{-1}(t)\dd t \stackrel{x \triangleq F_L^{-1}(t)}{=\joinrel=\joinrel=\joinrel=\joinrel=\joinrel=} \frac{1}{1-F_L(u)}\int_{u}^{+\infty} tf_L(t)\dd t,
\]
thus for $\nu > 1$
\begin{align*}
    \lim\limits_{\alpha \rightarrow 1^{-}} \frac{\ES_{\alpha}(L)}{\VaR_{\alpha}(L)} &= \lim\limits_{u \rightarrow +\infty} \frac{1}{u(1-F_L(u))}\int_{u}^{+\infty} tf_L(t)\dd t = \lim\limits_{u \rightarrow +\infty} \frac{\int_{u}^{+\infty} tf_L(t)\dd t}{cu^{-\nu+1}}\\
    &=\lim\limits_{u \rightarrow +\infty} \frac{-u f_L(u)}{c(-\nu+1)u^{-\nu}} =\lim\limits_{u \rightarrow +\infty} \frac{f_L(u)}{c(\nu-1)u^{-\nu-1}}\\
    &=\frac{\Gamma(\frac{\nu+1}{2})}{c(\nu-1)\sqrt{\nu\pi}\Gamma(\frac{\nu}{2})}\lim\limits_{u \rightarrow +\infty}\left(\frac{1}{u^2} + \frac{1}{\nu}\right)^{-\frac{\nu+1}{2}}\\
    &=\frac{1}{\nu-1}\nu^{-\frac{\nu-1}{2}}\nu^{\frac{\nu+1}{2}}\\
    &=\frac{\nu}{\nu-1} > 1.
\end{align*}
\qed

\HWproblem
Suppose the random variable $L \sim \Par(1/\xi)$, $0 < \xi < 1$. Show
\[
    \VaR_{\alpha}(L) = (1-\alpha)^{-\xi}, \quad \ES_{\alpha}(L) = \frac{\VaR_{\alpha}(L)}{1-\xi},
\]
and hence
\[
    \frac{\ES_{\alpha}(L)}{\VaR_{\alpha}(L)} = \frac{1}{1-\xi} > 1.
\]\\

\textbf{Proof:} For the Pareto distribution with shape \(1/\xi\) and scale \(1\),
\[
F_L(\ell)=1-\ell^{-1/\xi},\qquad \ell\ge1.
\]
For \(0<\alpha<1\), continuity and strict monotonicity give
\[
\alpha=F_L\bigl(\VaR_\alpha(L)\bigr)
=1-\VaR_\alpha(L)^{-1/\xi},
\]
so
\[
\VaR_\alpha(L)=(1-\alpha)^{-\xi}.
\]

And the expected shortfall
\[
\begin{aligned}
\ES_\alpha(L)
&=\frac{1}{1-\alpha}\int_\alpha^1\VaR_u(L)\,du=\frac{1}{1-\alpha}\int_\alpha^1(1-u)^{-\xi}\,du\\
&=\frac{1}{1-\alpha}\frac{(1-\alpha)^{1-\xi}}{1-\xi} = \frac{\VaR_\alpha(L)}{1-\xi},
\end{aligned}
\]
where the integral is finite because \(0<\xi<1\).

Consequently,
\[
\frac{\ES_\alpha(L)}{\VaR_\alpha(L)}
=\frac{1}{1-\xi}>1,
\]
since \(0<1-\xi<1\). \qed

\HWproblem
(Heavy-tailedness) Let random variables $L_1,L_2 \stackrel{\text{i.i.d.}}{\sim} \Par(1/2)$, which is, specifically, $\P(L_1 > x) = x^{-1/2}, x \ge 1$.

\HWsubproblem
Show that $\E(L_1) = \infty$\\

\textbf{Proof: } The probability density function of $L_1$ is
\[
    f_L(x) = [1-\P(L_1 > x)]' = (1-x^{-\frac12})' = \frac12 x^{-{\frac32}}, \quad x\ge 1.
\]
and $f_L(x) = 0$ if $x < 1$. Thus
\[
    \E(L_1) = \int_1^{+\infty} xf_L(x) \dd x =\int_1^{+\infty} \frac12 x^{-\frac12} \dd x = + \infty.
\]\qed

\HWsubproblem
Show that $L_1$ has a heavy tail, i.e.
\[
    \E(e^{tL_1}) = \infty \quad \text{for all } t>0.
\]\\

\textbf{Proof: }The moment generating function of $L_1$ at $t>0$ is
\begin{align*}
    \E(e^{tL_1}) = \int_1^{+\infty} e^{tx} f_L(x) \dd x = \int_1^{+\infty} e^{tx} \frac12x^{-\frac{3}{2}}\dd x = \frac12 \int_1^{+\infty} x^{-\frac32}e^{tx}\dd x,
\end{align*}
where for every fixed $t>0$ the integrand $x^{-\frac32}e^{tx}$ is positive and diverges to $+\infty$ as $x \rightarrow +\infty$: given any $M>0$ there is $x_0$ with $x^{-\frac32}e^{tx} \ge M$ for all $x \ge x_0$, so the integral is at least $\frac{M}{2}\int_{x_0}^{+\infty}\dd x = +\infty$. Hence $\E(e^{tL_1}) = +\infty$ for all $t>0$, i.e. $L_1$ has a heavy tail. \qed \\

\HWsubproblem
Show that $\P(L_1+L_2 > x) = 2\sqrt{x-1}/x, x\ge 2$.\\

\textbf{Proof: }Since $L_1$ and $L_2$ are independently and identically distributed, the joint PDF of $(L_1, L_2)$ is
\[
    f(u,v) = f_L(u)f_L(v) = 
    \begin{cases}
        \frac14u^{-\frac32}v^{-\frac32}, &u \ge 1 \text{ and } v \ge 1,\\
        0, &\text{otherwise}.
    \end{cases}
\] 
then
\[
    \P(L_1+L_2 > x) = 1- P(L_1 + L_2 \le x) = 1-\iint\limits_{\Omega} \frac14u^{-\frac32}v^{-\frac32} \dd u\mathrm{d}v
\]
where
\[
    \Omega = \{(u,v) \in \real^2: u \ge 1,v\ge 1,u+v \le x\}.
\]

Therefore,
\begin{align*}
        \P(L_1+L_2 > x) &= 1 - \frac14 \int_1^{x-1} \dd u \int_1^{x-u} u^{-\frac32}v^{-\frac32} \dd v = 1 - \frac{1}{2} \int_1^{x-1} u^{-\frac32}[1-(x-u)^{-\frac12}] \dd u \\
        & = 1 - \frac{1}{2} \int_1^{x-1}u^{-\frac32}\dd u + \frac{1}{2} \int_1^{x-1}u^{-\frac32}(x-u)^{-\frac12} \dd u \\
        & = \frac{1}{\sqrt{x-1}} + \frac12 \int_1^{x-1} u^{-\frac32}(x-u)^{-\frac12} \dd u.
\end{align*}
Let $u = x \sin^2 \theta$, $\theta \in [\alpha, \beta]$, where $x\sin^2\alpha = 1$ and $x\sin^2\beta = x-1$, which derive that
\[
    \alpha = \arctan \frac{1}{\sqrt{x-1}}, \quad \beta = \arctan \sqrt{x-1},
\]
then the integral
\begin{align*}
    \int_1^{x-1} u^{-\frac32}(x-u)^{-\frac12} \dd u &= \frac{2}{x} \int_{\alpha}^{\beta} \frac{\dd \theta}{\sin^2 \theta} = -\frac{2}{x} \cot \theta \Bigg|_{\alpha}^{\beta} \\
    &= -\frac{2}{x}(\frac{1}{\sqrt{x-1}}-\sqrt{x-1}) \\
    &=\frac{2(x-2)}{x\sqrt{x-1}}. 
\end{align*}
And therefore
\[
    \P(L_1+L_2 > x) = \frac{1}{\sqrt{x-1}} + \frac{x-2}{x\sqrt{x-1}} = \frac{2\sqrt{x-1}}{x},\quad x\ge2. 
\]\qed
    
\HWsubproblem
Show that for all $\alpha \in (0, 1)$,
\[
    \VaR_{\alpha}(L_1+L_2) > \VaR_{\alpha}(L_1) + \VaR_{\alpha}(L_2),
\]
that is, $\VaR_{\alpha}$ is superadditive.\\

\textbf{Proof: }Let $p = \VaR_{\alpha}(L_1+L_2)$, $q = \VaR_{\alpha}(L_1) = \VaR_{\alpha}(L_2)$. From the result of Problem 5,
\[
    q = \frac{1}{(1-\alpha)^2}
\]
(the same formula with $\xi = 2$, that is tail index $1/\xi = 1/2$; note that $\E(L_1) = \infty$, as shown above),
whereas, by the definition of value at risk, $\P(L_1+L_2 > p) = 2\sqrt{p-1}/p = 1- \alpha, \quad p \ge 2$, which gives
\[
    p = \frac{2(1+\sqrt{2\alpha-\alpha^2})}{(1-\alpha)^2},
\]
and the other solution
\[
    p' = \frac{2(1-\sqrt{2\alpha-\alpha^2})}{(1-\alpha)^2}
\]
has to be eliminated: the quantile must satisfy $p \ge 2$ (no smaller value is possible because $L_1+L_2 \ge 2$), whereas $p' < 2$ for every $\alpha \in (0,1)$. Indeed, with $s = 2\alpha-\alpha^2 \in (0,1]$,
\[
    p' \ge 2 \iff 1-\sqrt{s} \ge (1-\alpha)^2 \iff \sqrt{s} \le s \iff s = 1,
\]
so $p' \ge 2$ would force $\alpha = 1$, which is excluded by $\alpha \in (0,1)$.

Hence
\[
    p = \frac{2(1+\sqrt{2\alpha-\alpha^2})}{(1-\alpha)^2} > \frac{2}{(1-\alpha)^2} = 2q, \quad \forall \alpha \in (0,1),
\]
which is
\[
    \VaR_{\alpha} (L_1+L_2) > \VaR_{\alpha}(L_1) + \VaR_{\alpha}(L_2).
\]\qed

\HWproblem
(Skewness) In the market
\[
    L_i = 
    \begin{cases}
        -2, &\text{with probability 0.99}\\
        100, &\text{with probability 0.01}.
    \end{cases}
\]
There are two strategies (portfolios) for the investor to consider:\\
"Super Diversified": $\mathcal{P}_1 = \sum_{i=1}^{100} L_i$;\\
"Highly Concentrated": $\mathcal{P}_2 = 100L_1$.

\HWsubproblem
Show that the two strategies have the same expected loss.\\

\textbf{Proof: } the expected loss of a single bond
\[
    \E(L_i) = 0.99 \times (-2) + 0.01 \times 100 = -0.98,\quad \forall i \in \{1,2,\cdots, 100\}
\]
thus
\[
    \E(\mathcal{P}_1) = \E(\sum\limits_{i=1}^{100} L_i) = \sum\limits_{i=1}^{100} \E(L_i) = 100 \times (-0.98) = -98,
\]
\[
    \E(\mathcal{P}_2) = \E(100L_1) = 100\E(L_1) = 100 \times (-0.98) = -98.
\]
So the two strategies have the same expected loss. \qed

\HWsubproblem
Help the investor to analyze the potential risk of these two strategies by comparing
\[
    \VaR_{0.95}(\mathcal{P}_1) \text{  and  }\VaR_{0.95}(\mathcal{P}_2).
\]\\

Let $F_1(x), F_2(x)$ be the cumulative distribution functions of $\mathcal{P}_1$ and $\mathcal{P}_2$, respectively.
\[
    \P(\mathcal{P}_2 = -200) = 0.99, \quad \P(\mathcal{P}_2 = 10000) = 0.01,
\]
So
\[
    F_2(x) = \begin{cases}
        0, &x<-200, \\
        0.99, &-200 \le x < 10000,\\
        1, & x \ge 10000
    \end{cases}
\]
and
\[
    \VaR_{0.95}(\mathcal{P}_2) = \inf \{x\in \real: F_2(x) \ge 0.95\} = -200.
\]

Calculate the probability mass function (PMF) and the cumulative distribution function (CDF) for the first few possible values of $\mathcal{D}$.
\begin{table}[h]
    \centering
    \renewcommand{\arraystretch}{1.5} % stretch the rows to 1.5x
    \setlength{\tabcolsep}{20pt}     % widen the separation between columns to 20pt
    \begin{tabular}{r r r r}
        \hline
        $\mathcal{D}$ & $\mathcal{P}_1$ & PMF & CDF \\
        \hline
        0 & -200 & 0.366032 & 0.366032 \\
        1 & -98  & 0.369730 & 0.735762 \\
        2 & 4   & 0.184865 & 0.920627 \\
        3 & 106 & 0.060999 & 0.981626 \\
        4 & 208 & 0.014942 & 0.996568 \\
        \hline
    \end{tabular}
\end{table}
\\The CDF jumps past 0.95 at $\mathcal{D}= 3$, so 
\[
    \VaR_{0.95}(\mathcal{P}_1) = 106.
\]

So the "Highly Concentrated" portfolio has the smaller $\VaR$ — it looks safer at the 95\% level, even though it carries a 1\% chance of a 10,000 loss. This is the skewness of $\VaR$: it is not subadditive and it is blind to everything beyond the confidence quantile. The concentrated portfolio's whole tail risk lives in the 1\% mass above 0.95.

\end{document}
